\Needspace{6\baselineskip}
\begin{center}\hypertarget{native-251-article179}{}\pdfbookmark[2]{179.}{article-179}179.\end{center}
Iam ea quae ex hac lege sequuntur evolvemus. Sponte patet, ut productum
$\Omega=h^\mu\pi^{-\frac{1}{2}\mu}e^{-hh(vv+v'v'+v''v''+\mathord{.}\,\mathord{.})}$
fiat maximum, aggregatum $vv+v'v'+v''v''+\mathord{.}\,\mathord{.}$ minimum fieri
debere. \emph{Systema itaque maxime probabile valorum incognitarum $p$, $q$, $r$,
$s$ etc. id erit, in quo quadrata differentiarum inter functionum $V$, $V'$,
$V''$ etc. valores observatos et computatos summam minimam efficiunt,} siquidem
in omnibus observationibus idem praecisionis gradus praesumendus est.

Hocce principium, quod in omnibus applicationibus mathesis ad philosophiam
naturalem usum frequentissimum offert, ubique axiomatis loco eodem iure valere
debet, quo medium arithmeticum inter plures valores observatos eiusdem
quantitatis tamquam valor maxime probabilis adoptatur.

Ad observationes praecisionis \emph{inaequalis} principium nullo iam negotio
extendi potest. Scilicet si mensura praecisionis observationum, per quas inventum
est $V=M$, $V'=M'$, $V''=M''$ etc. resp. per $h$, $h'$, $h''$ etc. exprimitur,
i. e. si supponitur, errores his quantitatibus reciproce proportionales in istis
observationibus aeque facile committi potuisse, manifesto hoc idem erit, ac si
per observationes praecisionis aequalis (cuius mensura $=1$) valores functionum
$hV$, $h'V'$, $h''V''$ etc. immediate inventi essent $=hM$, $h'M'$, $h''M''$
etc.: \SourceLeaf{252}quamobrem systema maxime probabile valorum pro quantitatibus
$p$, $q$, $r$, $s$ etc. id erit, ubi aggregatum
$hhvv+h'h'v'v'+h''h''v''v''+\text{etc.}$, i. e. \emph{ubi summa quadratorum
differentiarum inter valores revera observatos et computatos per numeros qui
praecisionis gradum metiuntur multiplicatarum fit minimum.} Hoc pacto ne
necessarium quidem est, ut functiones $V$, $V'$, $V''$ etc. ad quantitates
homogeneas referantur, sed heterogeneas quoque (e. g. minuta secunda arcuum
et temporis) repraesentare poterunt, si modo rationem errorum, qui in singulis
aeque facile committi potuerunt, aestimare licet.
